\cotefile{157-3}
\scanpage{282}
\mspage{232}

une $W$-cofibration, $X'_0\unsure{-|\to}X'$ est dans $W$, et pour la $=$ raison $Y'_0-Y'$ l'est aussi. Donc $f'$ l'est aussi.

\uline{NB}\textsuperscript{\uline{bis}} C'est un argument formel, \inserted{(disons)} valable \inserted{(\unsure{sous} forme duale)} dans tout $(M,W)$ où on a un \unsure{théorème} de factorisation pour une flèche quelconque $f:X-Y$ en $X\xrightarrow{i_f}\overline X\xrightarrow{\alpha}Y$, avec $i_f\in W$ et $\alpha$ une $W$-fibration. Alors, si $X,Y$ sont $W$-fibrés sur $S$, \deleted{\unsure{toute|tout}} \unsure{une|e} $S$-flèche de $X$ dans $Y$ qui est dans $W$, est dans $W_S^{\unsure{\sim}}$. Et dualement, si \deleted{\unsure{toute}} flèche \inserted{$X-Y$} se factorise en $X'\longrightarrow\deleted{\unsure{X}}\widetilde Y\longrightarrow Y$, avec $X\longrightarrow\widetilde Y$ une $W$-cofibration et $\widetilde Y-Y$ \deleted{dans} $W$.

Il me vient ici une \inserted{(petite)} inquiétude : ai-je au moins prouvé que si
\[
S\longrightarrow X
\]
est une $W$-cofibration, alors pour toute flèche $S'\longrightarrow S$ dans $Cat$, le carré cocartésien dans $Cat$
\diagram{\begin{tikzcd}[ampersand replacement=\&]
 S \arrow[r,dash] \arrow[d] \& X \arrow[d] \\
 S' \arrow[r] \& X'
\end{tikzcd}}
est aussi $W$-\deleted{\unsure{Cat-cartésien|cocartésianique|cocartésien}} \inserted{cocartésien} \inserted{dans le sens ancien, spécial à $Cat$ (déf. 3 p. \unsure{131|191}) Pour le comparer à la théorie générale, cf prop 20 p. 262} ? \deleted{\unsure{Autrement|Effectivement},}

% archive sheet number: 282

\scanpage{283}
\mspage{233}

à quoi en vérité \unsure{serait-il|servait-il} bon ?! Pour le prouver, factorisons \inserted{encore} $S\unsure{-|\to}S'$ en
\diagram{\begin{tikzcd}[ampersand replacement=\&,column sep=large]
 S \arrow[r,hook,"{\substack{\text{imm.}\\\text{ouv. de}\\\text{Dwyer}}}"'] \& S'_0 \arrow[r,"{\in W}"'] \& S',
\end{tikzcd}}
d'où le diagramme de 2 carrés cocartésiens
\diagram{\begin{tikzcd}[ampersand replacement=\&]
 S \arrow[r,"{Cof_W}"] \arrow[d,hook,"{Dw}"'] \& X \arrow[d,hook,"{Dw}"] \\
 S'_0 \arrow[r,"{Cof_W}"] \arrow[d,"{\in W}"'] \& X'_0 \arrow[d,"{\in W}"] \\
 S' \arrow[r,"{Cof_W}"] \& X',
\end{tikzcd}}
où $X'_0-X'$ est $\in W$ d'après l'hypothèse sur $X$ \unsure{sous|sur} $S$. Il faut donc vérifier :

a) Un carré cocartésien construit sur une flèche qui est imm. de Dwyer, est \deleted{\unsure{homotopiq.}} $W$-cocartésien

b) Un carré commutatif \deleted{\unsure{dans|deux}} duquel deux côtés opposés sont dans $W$, est $W$-cocartésien.

c) Le composé de deux carrés $W$-cocartésiens est $W$-cocartésien.

Admettons provisoirement b) et c), et prouvons a), pour un carré coc.
\diagram{\begin{tikzcd}[ampersand replacement=\&,column sep=small,row sep=small]
 U \arrow[r,hook,"i"] \arrow[d,"p"'] \& X \arrow[d,"q"] \\
 U' \arrow[r,hook,"{i'}"] \& X'
\end{tikzcd}}
, où $\unsure{i|i'}$ est imm. \inserted{ouv.} de Dwyer. On procède \unsure{encore|alors} comme \unsure{tantôt|d'hab.}, en factorisant $p$ en $U\overset{\unsure{j\in|j}}{\hookrightarrow}U'_0\xrightarrow{p'}U'$, où $j$ est imm. ouv. de Dwyer, \inserted{$p'\in W$} \deleted{\unsure{On}}. d'où un diagramme en deux carrés cocartésiens accolés

% archive sheet number: 283

\scanpage{284}
\mspage{234}
\diagram{\begin{tikzcd}[ampersand replacement=\&]
 U \arrow[r,hook,"{i\in Dw}"] \arrow[d,hook,"{Dw}"'] \& X \arrow[d,hook] \\
 U'_0 \arrow[r,hook,"{Dw}"] \arrow[d,"{p'}"',"{\in W}"] \& X'_0 \arrow[d,"{q'\in W}"] \\
 U' \arrow[r,hook] \& X',
\end{tikzcd}}
où le premier carré est formé d'imm. \unsure{ouvertes|ouverte}. On sait \inserted{\unsure{alors|donc}} qu'il est \inserted{$W$-}cocartésien. \unsure{Ds|Dans} le deuxième, on a $p'\in W$, donc $q'\in W$ puisque $i$ est une imm. de Dwyer donc une $W$-cofibration. Donc il est égale\deleted{\unsure{ment}} $W$-cocartésien. On conclut alors grâce à \deleted{\unsure{b}}\unsure{c)} ci-dessus.

Prouvons b, pour
\diagram{\begin{tikzcd}[ampersand replacement=\&]
 X_0 \arrow[r,"{i_1}"] \arrow[d,"{i_2\in W}"'] \& X_1 \arrow[d,"{f_1\in W}"] \\
 X_2 \arrow[r,dash,"{f_2}"'] \& \deleted{\unsure{X}}Y
\end{tikzcd}}
, on l'insère dans
\diagram{\begin{tikzcd}[ampersand replacement=\&]
 X_0 \arrow[r,"{i_1}"] \arrow[d,"{i_2\in W}"'] \& X_1 \arrow[d,"{\unsure{j_1|i_1}}"] \arrow[ddr,bend left=20,"{f_1}"] \& \\
 X_2 \arrow[r,"{j_2}"'] \arrow[drr,bend right=24,"{f_2}"'] \& \unsure{\mathscr X|\mathcal X} \arrow[dr,"{\varphi}"] \& \\
 \& \& Y
\end{tikzcd}}
, où $\mathscr X=$
\diagram{\begin{tikzcd}[ampersand replacement=\&,row sep=tiny]
 \& \& X_1\\
 \displaystyle\int \& X_0 \arrow[ur] \arrow[dr] \&\\
 \& \& X_2
\end{tikzcd}}
. Comme $i_2\in W$, on sait qu'il en résulte $j_1\in W$, et comme $f_1\in W$, il en résulte $\varphi\in W$, \unsure{q.e.d.}

Reste à prouver c). Je vais le déduire d'une

\uline{Proposition} \uline{17}. Soit un carré commutatif
\diagram{\begin{tikzcd}[ampersand replacement=\&]
 \& X_0 \arrow[r,"{i_1}"] \arrow[d,"{i_2}"'] \& X_1 \arrow[d,"{f_1}"] \\
 Q={} \& X_2 \arrow[r,dash,"{f_2}"'] \& Y
\end{tikzcd}}
dans $Cat$, et $\mathscr X_Q$ la catégorie fibrée sur
\diagram{\begin{tikzcd}[ampersand replacement=\&]
 \& 0 \& 1 \arrow[l] \\
 Q^{\sim}(\simeq\Delta^1\times\Delta^1)={} \& 2 \arrow[u] \& 3 \arrow[l] \arrow[u]
\end{tikzcd}}
qui l'intègre. Soit $\mathscr X$ l'image inverse du fermé
\diagram{\begin{tikzcd}[ampersand replacement=\&,row sep=small,column sep=small]
 0 \& 1 \arrow[l] \\
 2 \arrow[u] \&
\end{tikzcd}}
% archive sheet number: 284

\scanpage{285}
\mspage{235}

dans $\unsure{\mathscr X|\mathcal X}_Q$, qui intègre le diagramme
\diagram{\begin{tikzcd}[ampersand replacement=\&,row sep=tiny]
 \& X_1\\
 X_0 \arrow[ur] \arrow[dr] \&\\
 \& X_2.
\end{tikzcd}}
Pour que $Q$ soit $W$-cocartésien, il f. et s. que l'inclusion
\[
i:\mathscr X\longrightarrow\mathscr X_Q
\]
soit \deleted{dans} $W$.

Comme $\mathscr X_Q$ est lisse sur $\widetilde Q$, et 3 est objet initial de $\widetilde Q$, il s'ensuit que l'inclusion de la fibre de \inserted{$\mathscr X_Q$ en le sommet} $3$ \deleted{de} $\unsure{\widetilde Q|Q}$ est dans $W$
\[
 Y\xrightarrow{i_Y}\deleted{\unsure{Q}}\mathscr X_Q\qquad\in W.
\]
Considérons \unsure{donc|alors} le diagramme
\diagram{\begin{tikzcd}[ampersand replacement=\&]
 \mathscr X \arrow[rr,"{\varphi_Q}"] \arrow[dr,"i"',""{name=left,coordinate,pos=.53}] \& \& Y \arrow[dl,"{i_Y\in W}", ""{name=right,coordinate,pos=.53}]\\
 \& \mathscr X_Q \arrow[from=right,to=left,dashed,"{\alpha}"'] \&
\end{tikzcd}}
Ce diagramme n'est pas commutatif, mais il y a une flèche de commutativité
\[
 \alpha:i_Y\varphi_Q\longrightarrow i.
\]
On sait que ceci implique que \deleted{$i\in W$}
\[
 i\in W \Longleftrightarrow (i_Y\varphi_Q)\in W
\]
(cf lemme p. 2), et comme $i_Y\in W$, \deleted{$i_Y\varphi_Q\in W\Longleftrightarrow$}
\[
 i_Y\varphi_Q\in W\Longleftrightarrow\varphi_Q\in W,
\]
d'où
\[
 i\in W\Longleftrightarrow\varphi_Q\in W\qquad \text{qed.}
\]
Nous pouvons maintenant prouver

% archive sheet number: 285

\scanpage{286}
\mspage{236}

\uline{Prop. 18} \inserted{(Moyennant $W(7\,\text{bis})$)} Considérons un diagramme commutatif
\diagram{\begin{tikzcd}[ampersand replacement=\&]
 \& X_0^{\deleted{\unsure{\prime}}} \arrow[rr] \arrow[dd] \& \& X_1^{\deleted{\unsure{\prime}}} \arrow[rr] \arrow[dd] \& \& X_2^{\deleted{\unsure{\prime}}} \arrow[dd,<->] \\
 (*) \& \& 1 \& \& 2 \& \\
 \& X'_0 \arrow[rr] \& \& X'_1 \arrow[rr] \& \& X'_2
\end{tikzcd}}
d'où un carré composé
\diagram{\begin{tikzcd}[ampersand replacement=\&,row sep=small,column sep=small]
 X_0 \arrow[rr] \arrow[dd] \& \& X_2 \arrow[dd] \\
 \& 3 \&\\
 X'_0 \arrow[rr] \& \& X'_2
\end{tikzcd}}
% slanted
\margin{\deleted{\unsure{Ça|Ce} devrait \unsure{être} vrai \unsure{dans} \unsure{toute|une} \unsure{catégorie} \illegible{2} \unsure{de|des} cofibrations}}

Supposons que le carré 1 soit $W$-cocartésien. Alors le carré 2 est $W$-cocartésien ssi le carré 3 l'est.

\uline{Dém} Soit
\diagram{\begin{tikzcd}[ampersand replacement=\&,column sep=small]
 I={} \& \bigg( \& 0 \arrow[d] \& 1 \arrow[l] \arrow[d] \& 2 \arrow[l] \arrow[d] \& \bigg)\\
 \& \& 0' \& 1' \arrow[l] \& 2' \arrow[l] \&
\end{tikzcd}}
et soit $\unsure{\mathscr X|\mathcal X}$ l'intégrale du diagramme $(*)$, catégorie fibrée sur $\mathscr X$. On désigne par $\mathscr X_{010'}$ la sous-catégorie pleine de $\mathscr X$ image inverse de la sous-catégorie pleine \deleted{$\unsure{\mathscr X|I}$}
\[
 \begin{aligned}
 I_{010'}&\subset I\\
 \rotatebox{90}{$=$}&\\
 \{0,1,0'\}&
 \end{aligned},
\]
et de $=$ pour toute autre partie \inserted{$A$} de \deleted{$Ob$} $I$ $Ob\,I=\{0,1,2,0',1',2'\}$, \deleted{\unsure{On pose}} \inserted{donnant une sous-catégorie pleine $\mathscr X_A$ de $\mathscr X$.} On pose
\diagram{\begin{tikzcd}[ampersand replacement=\&,column sep=small,row sep=small]
 U={} \& 0 \arrow[d] \& 1 \arrow[l] \& 2 \arrow[l] \\
 \& 0' \& \& 
\end{tikzcd}}
\diagram{\begin{tikzcd}[ampersand replacement=\&,column sep=small,row sep=small]
 ,\quad V={} \& 0 \arrow[d] \& 1 \arrow[l] \arrow[d] \& 2 \arrow[l] \\
 \& 0' \& 1' \arrow[l] \& ,
\end{tikzcd}}
d'où \deleted{$\unsure{\mathscr X}$}
\[
 \mathscr X_U=\mathscr X_{0120'},\qquad \mathscr X_V=\mathscr X_{0120'1'}.
\]
Les flèches entre les $\mathscr X_A$ sont les inclusions

% archive sheet number: 286

\scanpage{287}
\mspage{237}

déduites d'inclusions $A\subset A'$ \unsure{dans|ds} $\unsure{P|\mathcal P}(Ob\,I)$. Ainsi, le caractère $W$-cocartésien des carrés 1, 2, 3 s'\unsure{exprime|expriment} respectivement par la condition que les inclusions
\[
 \begin{aligned}
 \unsure{\mathscr X|\mathcal X}_{010'}&\xrightarrow{\varphi_1}\mathscr X_{010'1'}\\
 \mathscr X_{121'}&\xrightarrow{\varphi_2}\mathscr X_{121'2'}\\
 \mathscr X_{020'}&\xrightarrow{\varphi_3}\mathscr X_{020'2'}
 \end{aligned}
\]
soient dans $W$. Considérons \unsure{donc|alors} le diagramme de flèches d'inclusion
\diagram{\begin{tikzcd}[ampersand replacement=\&]
 \mathscr X_{020'} \arrow[r,"{\varphi_{\deleted{\unsure{2}}3}}"] \arrow[d,"{\alpha}"'] \& \mathscr X_{020'2'} \arrow[d,"{\beta}"] \\
 \mathscr X_U \arrow[r,"{\psi_3}"] \arrow[d,"{\deleted{\unsure{\varphi}}\psi_1}"'] \& \mathscr X \arrow[d,"{id}"] \\
 \mathscr X_V \arrow[r,"{\psi_2}"] \& \mathscr X\\
 \mathscr X_{\deleted{\unsure{020'}}121'} \arrow[r,"{\varphi_2}"] \arrow[u,"{\alpha'}"] \& \mathscr X_{\deleted{\unsure{020'2'}}121'2'} \arrow[u,"{\beta'}"']
\end{tikzcd}}
Je dis que dans ce diagramme on a
\[
 \alpha,\beta,\alpha',\beta'\ \in W \tag*{$(*)$}
\]
et
\[
 \varphi_1\in W\ \text{(i.e. carré 1 $W$-cocartésien)}
 \Longleftrightarrow\deleted{\unsure{\varphi}}\psi_1\in W. \tag*{$(**)$}
\]
Les relations $(*)$ impliquent déjà que
\[
 \begin{aligned}
 \varphi_2\in W&\Longleftrightarrow\psi_2\in W\\
 \varphi_3\in W&\Longleftrightarrow\psi_3\in W
 \end{aligned}
\]
et sous l'hyp. $W$-cocartésienne $(**)$, impliquant $\deleted{\unsure{\varphi}}\psi_1\in W$,

% archive sheet number: 287

\scanpage{288}
\mspage{238}

on trouve
\[
 \psi_2\in W\Longleftrightarrow\psi_3\in W,
\]
d'où
\[
 \varphi_2\in W\Longleftrightarrow\unsure{\beta_3|\varphi_3}\in W,
\]
ce qu'il fallait prouver.

Il reste à établir $(*)$ et $(**)$. Les relations $(*)$ proviennent de la remarque que si $i:A\longrightarrow A'$ est une inclusion de sous-catégories pleines de $I$, telle que $i$ soit $W$-\uline{asphérique}, \inserted{(i.e. $A_{/j}$ $W$-asphériques pour $j\in Ob\,A'\smallsetminus Ob\,A$),} alors $\unsure{\mathscr X|\mathcal X}_A\unsure{-|\to}\mathscr X_{A'}$ l'est aussi (chgt. de base lisse) donc est $\in W$. On doit vérifier la $W$-asphéricité pour les quatre inclusions
\[
 \begin{aligned}
 \{0,\deleted{\unsure{1}}2,0'\}&\longrightarrow U=\{0,\deleted{\unsure{0}}1,2,0'\}\\
 \{0,2,0',2'\}&\longrightarrow I\\
 \{1,2,1'\}&\longrightarrow V=\{0,1,2,0',1'\}\\
 \{1,2,1',2'\}&\longrightarrow I
 \end{aligned}
\]
\margin{$\Delta^0$\\[2pt] $\Delta^0\unsure{/|,}\Delta^1$.\\[2pt] $\Delta^1,\unsure{\Phi|Q}$.\\[2pt] $\Delta^1,\Delta^1\times\Delta^1$}

(NB j'ai marqué en marge les catégories qu'on obtient comme $A_{/j}$, \inserted{à iso. près,} dans les quatre cas envisagés -- elles ont \unsure{toutes} un objet initial, donc sont $W$-asphériques.)

Reste à prouver $(**)$. Soit
\[
 U'=\{0,1,0',\unsure{1'|1}\},\quad\text{donc}\quad
 V=U\cup U',\quad U\cap U'=\{0,1,0'\},
\]
\deleted{et \unsure{qui sont}} et $U,\unsure{W|U'}$ sont des \uline{fermés} de $V$ qui recouvrent $V$. Donc on aura de $=$
\[
 \mathscr X_V=\mathscr X_U\cup\mathscr X_{U'},\qquad
 \mathscr X_U\cap\mathscr X_{U'}=\mathscr X_{0,1,0'}
\]
% archive sheet number: 288

\scanpage{289}
\mspage{239}
\diagram{\begin{tikzcd}[ampersand replacement=\&]
 \unsure{\mathscr X|\mathcal X}_{U\cap U'} \arrow[r] \arrow[d,"{\varphi_1}"] \& \mathscr X_U \arrow[d,"{\psi_1}"]\\
 \mathscr X_{U'} \arrow[r] \& \mathscr X_V
\end{tikzcd}}
Donc par $W$(7 bis), (cas dual \deleted{\unsure{d}} des imm. fermées), on a $\varphi_1\in W\Longleftrightarrow\psi_1\in W$, qed.

Notons pour mémoire le corollaire qui nous avait motivé

\uline{Corollaire} Supposons $W$(7 bis) \inserted{(comme dans \unsure{tte} cette section)}. Soit un carré cocartésien dans $Cat$
\diagram{\begin{tikzcd}[ampersand replacement=\&]
 X_0 \arrow[r,"{i_1}"] \arrow[d,"{i_2}"'] \& X_1 \arrow[d]\\
 X_2 \arrow[r] \& X_3,
\end{tikzcd}}
avec $i_1$ (ou $i_2$) une $W$-cofibration (p. ex. une immersion de Dwyer). Alors ce carré est $W$-cocartésien.

Cette digression (p. 232 ff) n'est pas \unsure{nécessaire} d'ailleurs pour formuler dès maintenant la conclusion de \unsure{nos|mes} cogitations sur les cofibres $W$-homotopiques : \uline{Il y a une belle théorie dans} $Cat$, pourvu seulement que $W$ satisfasse à $W$(7 bis). En effet \unsure{tout} résulte du fait que \unsure{toute} flèche \inserted{$f:X\longrightarrow Y$} dans $Cat$ se factorise en
\[
 X\xrightarrow{i_f\deleted{\unsure{p}}}\widetilde Y\xrightarrow{\unsure{p_f|\rho_f}}Y
\]
avec $i_f$ une $W$-cofibration (et $=$ une imm

% archive sheet number: 289

\scanpage{290}
\mspage{240}

ersion ouverte de Dwyer), et $p_f\in W$.

7. \uline{Retour sur l'axiomatique de la fibre $W$-homotopique}.

Je reviens sur \unsure{les réflexions|la réflexion} du n\textsuperscript{o} 4. Finalement, pour
\[
 (M,W\subset Fl\,M)
\]
\unsure{donnés}, avec $W$ faiblement saturé, il apparaît qu'une excellente théorie des fibres $W$-homotopiques, satisfaisant aux conditions \textcircled{a} \textcircled{b} \textcircled{c} du n\textsuperscript{o} 4 (page 215-218), existe \unsure{déjà|dès|sitôt} que l'axiome \deleted{\unsure{anodin|ci-dessous}} d'apparence anodine suivant est satisfait (en plus de l'existence \inserted{de l'objet final et} des produits fibrés)

% slanted
\margin{et encore ! Il me semble qu'il suffit de supposer que les $X\to Y$ qui sont $W$-\unsure{fibrés} sont ``\unsure{quarrables}''.}

\uline{Factorisation gauche} : Toute flèche \inserted{$f:\unsure{\alpha|X}\unsure{-|\to}Y$} de $M$ se factorise en
\[
 \begin{aligned}
 X&\xrightarrow{i_f}\overline X\xrightarrow{p_f}Y,\\
 i_f&\in W,\quad p_f\in Fib_W.
 \end{aligned}
\]
% slanted
\margin{\uline{NB} L'axiome de \deleted{\unsure{suffisant|factorisation}} factorisation \unsure{à gauche|gauche} n'est suffisant pour l'existence d'une bonne théorie des fibres homotopiques \inserted{(\unsure{semble-t-il|malentendu})}, que dans le cas où tout $X$ de $M$ est $W$-fibrant, i.e. est dans $\underline{Fib}_W\,e$.}

% slanted
\margin{Sous ces conditions $M$ muni de $(W,F_W)$ est une \unsure{catégorie|catég.} à \unsure{fibrations|cofibrations} de K.S. \deleted{\unsure{Br.}} Brown -- \unsure{Au dessous|Au dessus} \unsure{de|cf.} ``tout est OK'' mais il faudrait vérifier a) b) c) ci-dessus (p. 215, 218) je ne suis pas sûr de l'avoir fait}

L'axiome dual donne lieu à une bonne théorie des cofibres $W$-homotopiques :

\uline{Factorisation droite} : Toute flèche $f:X-Y$ dans $M$ se factorise en
\[
 X\xrightarrow{j_f}\widetilde Y\xrightarrow{q_f}Y,
\]
avec $q_f\in W$, $j_f\in Cofib_W$.

% archive sheet number: 290

\scanpage{291}
\mspage{241}

Je voudrais à présent développer quelque peu le \unsure{formalisme|fonctorisme} des fibres $W$-homotopiques, sous l'hypothèse \inserted{(disons)} que \inserted{$W$ satisfait} l'axiome de factorisation gauche.

Considérons un diagramme cartésien
\diagram{\begin{tikzcd}[ampersand replacement=\&]
 X_1 \arrow[d,Rightarrow,"{p_1}"'] \& \deleted{\unsure{X_3}}Y \arrow[l,"{q_2}"'] \arrow[d,"{q_1}"]\\
 \deleted{\unsure{X_0}}S \& X_2 \arrow[l,"{p_2}"]
\end{tikzcd}}
dans $M$. Si $p_1$ est dans $Fib_W$, alors $Y$ \deleted{\unsure{peut}} \inserted{sur $X_2$} peut s'interpréter comme la fibre \inserted{$W$-}homotopique de $X_1/S$ en $X_2/S$. Plus précisément, il en est ainsi de la classe $ho^{\unsure{lc|le}}_{W,X_2}(Y)$,
\[
 ho^{lc}_{\deleted{\unsure{W}}X_2,\deleted{\unsure{X_2}}W}
 (Y/\deleted{\unsure{S}}X_2)\simeq p_2^*\bigl(ho^{lc}_{S,W}(X_{\unsure{1|a}}/S)\bigr)
 \tag*{$(*)$}
\]
Mais en tant qu'objet sur $X_1$, $Y$ peut s'interpréter également comme fibre \inserted{$W$-}homotopique de $X_2/S$ en $X_1/S$ :
\[
 ho^{lc}_{X_1,W}(Y/X_1)\simeq p_1^*\bigl(ho^{lc}_{S,W}(X_2/S).
 \tag*{$(**)$}
\]
A dire vrai, il faudrait le prouver, Pour ceci, considérons \unsure{une|leur|la} factorisation
\[
 X_2\xrightarrow{i_2}\overline X_2\xrightarrow{\overline p_2}S
 \qquad i_2\in W,\ \overline p_2\in Fib_W,
\]
de sorte qu'on trouve un diagramme de carrés cartésiens
\diagram{\begin{tikzcd}[ampersand replacement=\&]
 X_1 \arrow[d,Rightarrow] \& \overline Y \arrow[l,"{\overline q_2}"'] \arrow[d] \& Y \arrow[l,"{j_2}"',"W"] \arrow[d]\\
 S \& \overline X_2 \arrow[l,"{\overline p_2}"'] \& X_2 \arrow[l,"{i_2}"',"W"]
\end{tikzcd}}
% archive sheet number: 291

\scanpage{292}
\mspage{242}

Par définition \inserted{(de $p_1^*$)} \deleted{\unsure{de}}, on a
\[
 p_1^*\bigl(ho^{\unsure{lc|le}}_{W,S}(X_{\unsure{2|1}}/S)\bigr)
 \simeq \deleted{\unsure{p}}ho^{lc}_{W,X_1}(\overline Y/X_1),
\]
d'autre part, comme $X_1\unsure{-|\to}S$ est \unsure{une|un} $Fib_W$, \deleted{$j_2$} et $i_2\in W$, on a $j_2\in W$, donc
\[
 ho^{lc}_{W,X_1}(Y/X_1)\xrightarrow{\sim}ho^{lc}_{W,X_1}(\overline Y/X_1),
\]
d'où la formule $(**)$.

Ainsi, l'objet $Y$ a deux interprétations homotopiques différentes, suivant qu'on le regarde comme étant sur $X_1$, \unsure{ou|soit} sur $X_2$. On aimerait donc une interprétation qui englobe les deux à la fois. \deleted{\unsure{Mon propos|Ma propos.}} \inserted{\unsure{On va}} \deleted{\unsure{est de}} considérer le diagramme
\diagram{\begin{tikzcd}[ampersand replacement=\&]
 X_1\arrow[d] \&\\
 S \& X_2\arrow[l]
\end{tikzcd}}
de type $\Psi^{\unsure{o|0}}$, où
\diagram{\begin{tikzcd}[ampersand replacement=\&,row sep=small,column sep=small]
 \& 1 \&\\
 \Psi={} \& 0\arrow[u]\arrow[r] \& 2
\end{tikzcd}}
ou
\diagram{\begin{tikzcd}[ampersand replacement=\&,row sep=tiny]
 \& 1\\
 0\arrow[ur]\arrow[dr] \&\\
 \& 2,
\end{tikzcd}}
dont $Y$ se déduit par produit fibré, en complétant ce diagramme en un diagramme de type $\widetilde\Psi^{o}$, où
\diagram{\begin{tikzcd}[ampersand replacement=\&,row sep=tiny,column sep=small]
 \& \& 1 \arrow[dr] \& \&\\
 \deleted{\unsure{Q=}}\ \widetilde\Psi={} \& 0\arrow[ur]\arrow[dr] \& \& 3 \& (\simeq\Delta^1\times\Delta^1)\\
 \& \& 2\arrow[ur] \& \&
\end{tikzcd}}
% archive sheet number: 292

\scanpage{293}
\mspage{243}

Pour toute catégorie $I$, soit
\[
 Ho_{M,W}(I)\ \text{ou}\ Ho_W(I)\ \overset{\mathrm{def}}{=}\ M(I)\,W(I)^{-1},
\]
où
\[
 \left\{\begin{aligned}
 M(I)&=\underline{Hom}(I^{\unsure{o|0}},M)\\
 W(I)&\subset Fl(M(I))\quad\text{est formé des }u:F\unsure{-|\to}G\\
 &\quad\text{telles que }u_i:F(i)\to G(i)\text{ dans }W\\
 &\quad\text{pour tout }i\text{ dans }Ob\,I.
 \end{aligned}\right.
\]
Ainsi, le diagramme
\diagram{\begin{tikzcd}[ampersand replacement=\&,row sep=tiny,column sep=small]
 \& X_1\arrow[dl]\\ S\&\\ \& X_2\arrow[ul]
\end{tikzcd}}
définit un objet \inserted{$\xi$} de $M(\Psi)$, donc de $Ho_W(\Psi)$, et le \deleted{\unsure{di}} diagramme
\diagram{\begin{tikzcd}[ampersand replacement=\&,row sep=tiny,column sep=small]
 \&X_1\arrow[dl]\&\\ S\&\&Y\arrow[ul]\arrow[dl]\\ \&X_2\arrow[ul]\&
\end{tikzcd}}
peut de $=$ s'interpréter comme un \deleted{objet} \inserted{$\widetilde\xi$} de $Ho_W(\widetilde\Psi)$, \unsure{dont|d'où} la restriction à $\Psi$ est $\overset{\deleted{\unsure{\sim}}}{\xi}$.

% slanted
\margin{\uline{NB} Cette construction, \unsure{qui se|que je} fait ici assez \unsure{lourdement}, dérive \deleted{\unsure{de façon}} \unsure{immédiatemt|immédiatement|immédiate}, et \unsure{ainsi|aussi} \unsure{ses} principales propriétés, de \unsure{l'énoncé} que la cat. à fibrations $(M,\overline W,Fib_W)$ définit un \uline{dérivateur} sur le \deleted{\unsure{domaine}} \unsure{domaine} des ens. ordonnés finis, ce qui \deleted{\unsure{se|en}} \unsure{prouve|pourra} en effet \unsure{assez|s'app.} simplement (sans \unsure{pas} parler du mal que je me donne ici).}

Ceci posé, j'aimerais définir un \unsure{foncteur}
\[
 i_*:Ho_W(\Psi)\longrightarrow Ho_W(\widetilde\Psi)
\]
\[
 \unsure{v|u|i}:\Psi\longrightarrow\Psi^{\sim}
\]
l'inclusion, qui est une imm. ouverte), lequel doit correspondre, dans le cas d'un objet
\diagram{\begin{tikzcd}[ampersand replacement=\&,row sep=tiny,column sep=small]
 \&X_1\arrow[dl]\\ S\&\\ \&X_2\arrow[ul]
\end{tikzcd}}
de $Ho_W(\Psi)$ avec $p_1\in Fib_W$, à l'opération de produit fibré.

% archive sheet number: 293

\scanpage{294}
\mspage{244}

La construction de ce foncteur $i_*$ devrait être essentiellement triviale, avec ce dont nous disposons. Nous allons \deleted{\unsure{le}} définir d'abord, \deleted{sur la sous-catégorie} pour \uline{$S$ fixé}, un foncteur canonique
\[
 Ho^{\unsure{lc|le}}_{W\deleted{\unsure{,}}}(S)\times Ho^{lc}_{W}(S)
 \longrightarrow Ho_{W,\deleted{\unsure{S'}}S}(\Phi) \tag*{$(*)$}
\]
où le deuxième membre est défini, par analogie avec les $Ho_W(I)$, par
\[
 Ho_{W,S}(I)\overset{\mathrm{def}}{=}
 \bigl(M_{/S}(I)\bigr)\bigl(W_S(I)\bigr)^{-1}
\]
-- i.e. c'est la $=$ définition, sauf que $(M,W)$ est remplacé par $(M_{/S},W_S)$. Quant à
\diagram{\begin{tikzcd}[ampersand replacement=\&,row sep=tiny,column sep=small]
 \&1\arrow[dr,"{\deleted{\unsure{\nwarrow}}}"{near start}]\&\\ \Phi={} \&\&3,\\ \&2\arrow[ur,"{\deleted{\unsure{\swarrow}}}"{near start}]\&
\end{tikzcd}}
on le considère comme une sous-catégorie fermée de $\Psi^{\sim}$, qui s'en déduit en lui rajoutant un objet initial $0$ (tout comme $\Psi^{\sim}$ se déduit de $\Psi$ en lui rajoutant un objet final $3$). D'autre part, on a un foncteur tautologique
\[
 Ho_{W,S}(\Phi)\longrightarrow Ho_W(\deleted{\unsure{\Phi}}\unsure{\Psi^{\sim}|\widetilde\Psi}),
 \tag*{$(*')$}
\]
qui sur les objets consiste à compléter un diagramme
\diagram{\begin{tikzcd}[ampersand replacement=\&,row sep=tiny,column sep=small]
 X_1\&\\ \&Y\arrow[ul]\arrow[dl]\\ X_2\&
\end{tikzcd}}
de $\deleted{\unsure{Cat}}M_{/S}$ en le diagramme carré
\diagram{\begin{tikzcd}[ampersand replacement=\&,row sep=tiny,column sep=small]
 \&X_1\arrow[dl]\&\\ S\&\&Y\arrow[ul]\arrow[dl]\\ \&X_2\arrow[ul]\&
\end{tikzcd}}
de $M$.

% archive sheet number: 294

\scanpage{295}
\mspage{245}

Le composé de $(*)$ et de $(*')$ \unsure{donnera|sera} \inserted{\unsure{donc}} un foncteur
\[
 Ho^{\unsure{lc|le}}_W(S)\times Ho^{lc}_W(S)\longrightarrow Ho_W(\Psi^{\sim}),
 \tag*{$(**)$}
\]
très proche déjà du foncteur $i_*$ à construire. Pour construire $(*)$, on interprète (à équiv. de catégories près) le premier membre comme
\[
 \bigl(\underline{Fib}_W\,S\deleted{\unsure{)}}\times
 \underline{Fib}_W\,S\bigr)(W_S\times W_S)^{-1}
\]
donc il faut construire un foncteur
\[
 \underline{Fib}_W\,S\times\underline{Fib}_W\,S\longrightarrow Ho_{W,S}(\Phi),
\]
qui transforme flèches dans $W_S\times W_S$ en iso. On prendra le foncteur
\diagram{\begin{tikzcd}[ampersand replacement=\&,row sep=small,column sep=small]
 (X_1,X_2)\arrow[r,mapsto] \& \bigg( \& X_1 \& X_1\times_S X_2\arrow[l]\arrow[d] \&\bigg)\\
 \&\&\&X_2\&,
\end{tikzcd}}
et il faut seulement vérifier que si
\[
 X_1\xrightarrow{\alpha_1}X'_1,\quad X_2\overset{\alpha_2}{\unsure{-|\to}}X'_2
\]
sont deux flèches de $\underline{Fib}_W\,S$, \deleted{\unsure{alors}}
\[
 \deleted{\unsure{\alpha|X}}X_1\times_S X_2\xrightarrow{\alpha=\alpha_1\times_S\alpha_2}X'_1\times_S X'_2,
\]
alors
\[
 \alpha_1\in W,\ \alpha_2\in W\Longrightarrow\alpha\in W,
\]
ce qui est immédiat sur la factorisation
\diagram{\begin{tikzcd}[ampersand replacement=\&,column sep=large]
 X_1\times_S X_2 \arrow[r,"{\substack{\rotatebox{-90}{$\in$}\\ W\\ \text{car }X_2\in\underline{Fib}_W\,S}}"'] \&
 X'_1\times_S X_2 \arrow[r,"{\substack{\rotatebox{-90}{$\in$}\\ W\\ \text{car }X'_1\in\underline{Fib}_W\,S}}"'] \&
 X'_1\times_S X'_2 .
\end{tikzcd}}
% archive sheet number: 295

\scanpage{296}
\mspage{246}

On a prouvé la

\uline{Proposition 1}8 Soit $S$ dans $M$. Il existe, à isom unique près, un foncteur et un seul
\[
 Ho^{\unsure{lc|le}}_{W\deleted{\unsure{,}}}(S)\times Ho^{lc}_W(S)\xrightarrow{\pi}Ho_{W,S}(\Phi)
\]
\deleted{\unsure{dont le composé avec}} \inserted{donnant lieu à un carré commutatif} à
\[
 \deleted{\underline{Fib}_W\,S\times\underline{Fib}_W\,S\longrightarrow Ho^{lc}_W(S)\times Ho^{lc}_W(S)}
\]
isom près
\diagram{\begin{tikzcd}[ampersand replacement=\&,column sep=2pt,row sep=tiny,font=\scriptsize]
 \& \& X_1 \& \\
 \& (X_1,X_2)\arrow[r,mapsto] \& {} \& X_1\times_S X_2\arrow[ul]\arrow[dl] \\
 \& \& X_2 \& \\
 \underline{Fib}_W\,S\times\underline{Fib}_W\,S\arrow[rrr]\arrow[d] \& \& \& \deleted{\unsure{\underline{Fib}_W}}M_{/S}(\Phi)\arrow[d]\\
 Ho^{lc}_W(S)\times Ho^{lc}_W(S)\arrow[rrr] \& \& \& Ho_{W,S}(\Phi)\\
 \bigl(=(M_{/S})(W_S)^{-1}\times(M_{/S})(W_S)^{-1} \& \& \& \bigl(=M_{/S}(\Phi)(W_S(\Phi))^{-1}\bigr)
\end{tikzcd}}
\deleted{\unsure{Consid.}} Pour $X_1,X_2$ dans $M_{/S}$, \deleted{\unsure{leur}} le diagramme
\[
 \pi(X_1,X_2)\overset{\mathrm{def}}{=}\pi\bigl(ho^{lc}_{W,S}(X_1),ho^{lc}_{W,S}(X_2)\bigr)
\]
se ``calcule'' ainsi à iso. can. près : on choisit
\[
 X_1\overset{\alpha_1}{\unsure{-|\to}}\overline X_1,\quad
 X_2\overset{\alpha_2}{-}\overline X_2\qquad\text{flèches dans }M_{/S}
\]
avec
\[
 \alpha_1,\alpha_2\in W_S,\qquad\overline X_1,\overline X_2\in\mathrm{Ob}\,\underline{Fib}_W\,S,
\]
ce qui définit
\diagram{\begin{tikzcd}[ampersand replacement=\&,column sep=small,row sep=small]
 \pi(\alpha_1,\alpha_2):\quad \pi(X_1,X_2)\arrow[r,"\sim"] \&
 \pi(\overline X_1,\overline X_2)=ho^{\unsure{\xi|\zeta}}_{\unsure{S,M,\Phi|S,W,\Phi}}\bigg( \&
 \overline X_1 \& \overline X_1\times_S\overline X_2\arrow[l]\arrow[d] \& \bigg)\\
 \& \& \&\overline X_2\& .
\end{tikzcd}}
% archive sheet number: 296

